\chapter{Séries de Laurent e o teorema dos resíduos}
\label{ch:b3:residues}

O que acontece com uma \omterm{def:b3:holomorphic:holo}{função holomorfa} perto de um ponto em que ela
\emph{não} está definida? A resposta é uma tricotomia \omterm{def:b3:complete:complete}{completa} ---
ponto removível, polo ou singularidade essencial --- que se lê em uma
série de potências de duas pontas, a expansão de Laurent. Um
coeficiente dessa expansão, o \emph{\omterm{def:b3:residues:singularities}{resíduo}}, controla toda \omterm{def:b3:holomorphic:contour}{integral de
contorno} em torno da singularidade: o teorema dos \omterm{def:b3:residues:singularities}{resíduos} converte
integrais definidas difíceis em álgebra finita, conta zeros de funções
(princípio do argumento, Rouché) e demonstra o teorema da aplicação
aberta. Antes disso, elevamos o teorema de Cauchy dos domínios
estrelados à sua forma definitiva, sem homologia --- o elegante
argumento de Dixon ---, de modo que todos os contornos de índice zero em
torno do complementar fiquem disponíveis.

\section{O teorema de Cauchy global}

Um \emph{ciclo}\index{ciclo} $\Gamma$ é uma soma formal finita de
\omterm{def:b3:topology:pathconnected}{caminhos} fechados $\gamma_1, \dots, \gamma_m$; as integrais e os índices ao longo de
$\Gamma$ são as somas correspondentes, e $\operatorname{im}\Gamma = \bigcup\operatorname{im}\gamma_j$.

\begin{theorem}[Cauchy, forma global]
\label{thm:b3:residues:globalcauchy}
Sejam $\Omega \subseteq \C$ aberto, $f \in \mathcal H(\Omega)$ e $\Gamma$ um ciclo em $\Omega$ tal que
\[
\operatorname{Ind}_\Gamma(w) = 0
\qquad\text{para todo } w \notin \Omega .
\]
Então, para todo $z \in \Omega\setminus\operatorname{im}\Gamma$,
\[
\frac1{2\iu\pi}\int_\Gamma\frac{f(w)}{w - z}\,\dd w =
\operatorname{Ind}_\Gamma(z)\,f(z),
\qquad\text{e}\qquad
\int_\Gamma f(w)\,\dd w = 0 .
\]
\index{teorema de Cauchy (global)}
\end{theorem}

\begin{proof}[Demonstração (Dixon)]
Defina $g \colon \Omega\times\Omega \to \C$ por
\[
g(z, w) = \begin{cases}
\dfrac{f(w) - f(z)}{w - z} & w \neq z,\\[4pt]
f'(z) & w = z .
\end{cases}
\]
\emph{$g$ é \omterm{def:b3:topology:continuity}{contínua}}: fora da diagonal, é claro. Perto de um ponto
diagonal $(a, a)$, expanda $f$ em série de potências em $a$
(\cref{thm:b3:holomorphic:analytic}): $f(w) - f(z) =
\sum_{n\geq1}c_n\bigl((w-a)^n - (z-a)^n\bigr)$, e dividindo cada termo por $w - z$ (fatoração
de $u^n - v^n$) obtém-se, para $z, w \in D(a, r)$,
\[
g(z, w) =
\sum_{n\geq1}c_n\sum_{j=0}^{n-1}(w-a)^{\,j}(z-a)^{\,n-1-j},
\]
válido também na diagonal (cada soma interna vira $n(z-a)^{n-1}$, somando
$f'(z)$). Para $r$ pequeno, a série converge uniformemente em $D(a,r)^2$
($\abs{\text{termo}} \leq n\abs{c_n}r^{n-1}$, somável dentro do raio): a soma é \omterm{def:b3:topology:continuity}{contínua}.

Ponha $h(z) = \frac1{2\iu\pi}\int_\Gamma g(z, w)\,\dd w$ em $\Omega$: \omterm{def:b3:topology:continuity}{contínua} (\omterm{def:b3:topology:continuity}{continuidade} uniforme de $g$ nos
\omterm{def:b3:topology:compact}{compactos}) e \omterm{def:b3:holomorphic:holo}{holomorfa} --- por Morera (critério do~\cref{thm:b3:holomorphic:weierstrassconv}): para
um triângulo $T \subseteq \Omega$, Fubini dá $\int_{\partial
T}h = \frac1{2\iu\pi}\int_\Gamma\bigl(\int_{\partial T}g(z,
w)\dd z\bigr)\dd w = 0$, anulando-se a integral interna porque
$z \mapsto g(z, w)$ é \omterm{def:b3:holomorphic:holo}{holomorfa} em $\Omega$ (em $z = w$ a singularidade é
removível: $g$ é ali \omterm{def:b3:topology:continuity}{contínua} e \omterm{def:b3:holomorphic:holo}{holomorfa} nos demais pontos --- o
argumento de extensão da demonstração do~\cref{thm:b3:holomorphic:formula}).

No aberto $\Omega' = \{z \notin \operatorname{im}\Gamma
: \operatorname{Ind}_\Gamma(z) = 0\}$, defina $h_1(z) =
\frac1{2\iu\pi}\int_\Gamma\frac{f(w)}{w - z}\dd w$: \omterm{def:b3:holomorphic:holo}{holomorfa} em $\Omega'$ (Morera ou derivação sob
a integral). Para $z \in \Omega\cap\Omega'$:
\[
h(z) = \frac1{2\iu\pi}\int_\Gamma\frac{f(w)}{w-z}\dd w
- f(z)\operatorname{Ind}_\Gamma(z) = h_1(z) .
\]
Por hipótese, $\Omega\cup\Omega' = \C$ ($w \notin \Omega
\Rightarrow \operatorname{Ind}_\Gamma(w) = 0$), de modo que $h$ e $h_1$ se colam em uma
função inteira $H$. Como a componente ilimitada do complementar de
$\operatorname{im}\Gamma$ está em $\Omega'$ e $h_1(z) \to 0$ quando $\abs z \to \infty$ (cota ML), $H$ é limitada e
tende a $0$: Liouville (\cref{cor:b3:holomorphic:liouville}) dá $H \equiv 0$. Assim $h \equiv 0$ em
$\Omega$, o que é a fórmula integral. Aplicando-a, para $a \in
\Omega\setminus\operatorname{im}\Gamma$ fixado,
a $\tilde f(w) =
(w - a)f(w)$ em $z = a$:
\[
\frac1{2\iu\pi}\int_\Gamma f(w)\dd w =
\frac1{2\iu\pi}\int_\Gamma \frac{\tilde f(w)}{w - a}\dd w =
\operatorname{Ind}_\Gamma(a)\,\tilde f(a) = 0 .
\]
\end{proof}

\section{Séries de Laurent e singularidades isoladas}

\begin{theorem}[Expansão de Laurent]\label{thm:b3:residues:laurent}
Seja $f$ \omterm{def:b3:holomorphic:holo}{holomorfa} no anel $A = \{r < \abs{z - a} <
R\}$ ($0 \leq r < R \leq \infty$). Então\index{série de Laurent}
\[
f(z) = \sum_{n\in\Z}c_n\,(z - a)^n
\qquad\text{em } A,
\qquad
c_n = \frac1{2\iu\pi}\int_{C_\rho}\frac{f(w)}{(w -
a)^{n+1}}\,\dd w
\]
para qualquer $r < \rho < R$ (independente de $\rho$), convergindo as duas
semisséries normalmente em subanéis \omterm{def:b3:topology:compact}{compactos}. A expansão é única.
\end{theorem}

\begin{proof}
Fixe $r < \rho_1 < \abs{z - a} < \rho_2 < R$ e seja $\Gamma =
C_{\rho_2} - C_{\rho_1}$ (o externo no sentido anti-horário, o interno no
horário): um ciclo em $A$ com $\operatorname{Ind}_\Gamma(w) = 0$ para todo $w \notin A$ (pontos dentro do
disco pequeno: $1 - 1 = 0$; fora do grande: $0 - 0$). Pelo~\cref{thm:b3:residues:globalcauchy},
$\operatorname{Ind}_\Gamma(z) = 1 - 0 = 1$ dá
\[
f(z) = \frac1{2\iu\pi}\int_{C_{\rho_2}}\frac{f(w)}{w - z}\dd
w - \frac1{2\iu\pi}\int_{C_{\rho_1}}\frac{f(w)}{w - z}\dd w .
\]
Expanda o primeiro núcleo como no~\cref{thm:b3:holomorphic:analytic} (potências de $\frac{z -
a}{w - a}$, de
módulo $< 1$): a parte não negativa $\sum_{n\geq0}c_n(z-a)^n$. No segundo, expanda ao
contrário: $\frac{-1}{w - z} = \frac1{(z-a)(1 - \frac{w - a}{z -
a})} = \sum_{m\geq0}\frac{(w-a)^m}{(z - a)^{m+1}}$, normalmente convergente em $C_{\rho_1}$: a parte negativa
$\sum_{n\leq-1}c_n(z-a)^n$, com os coeficientes enunciados (índice $n = -m-1$). Independência de
$\rho$: as integrais dos coeficientes em $C_{\rho}$ e em $C_{\rho'}$ diferem por
$\int_\Gamma$ sobre um ciclo de índice nulo da \omterm{def:b3:holomorphic:holo}{função holomorfa} $\frac{f(w)}{(w-a)^{n+1}}$ em $A$:
zero, de novo pelo~\cref{thm:b3:residues:globalcauchy}. Unicidade: integre $\sum c_n(z-a)^n$ contra $(z - a)^{-m-1}$
sobre $C_\rho$ termo a termo (convergência normal): só $n = m$
sobrevive.
\end{proof}

\begin{definition}\label{def:b3:residues:singularities}
Se $f$ é \omterm{def:b3:holomorphic:holo}{holomorfa} em um disco perfurado $D(a, R)\setminus
\{a\}$, expanda por Laurent
($r = 0$). Três casos exclusivos:\index{singularidade isolada}
\begin{itemize}
  \item todos os $c_n = 0$ para $n < 0$: \emph{\omterm{thm:b3:residues:riemanncw}{singularidade removível}}
        (a série não negativa estende $f$ \omterm{def:b3:holomorphic:holo}{holomorficamente} a
        $a$);
  \item $c_n \neq 0$ para um número finito, e ao menos um, de $n < 0$: um
        \emph{polo} de ordem $m = -\min\{n : c_n \neq 0\}$; equivalentemente, $f = g/(z-a)^m$ com $g$
        \omterm{def:b3:holomorphic:holo}{holomorfa} e $g(a)
        \neq 0$; equivalentemente, $\abs{f(z)} \to \infty$ quando $z\to
        a$;
  \item infinitos $c_n \neq 0$ negativos: \emph{singularidade essencial}.
\end{itemize}
O \emph{resíduo}\index{resíduo} é $\operatorname{Res}(f, a) = c_{-1}$. Uma \omterm{def:b3:holomorphic:holo}{função holomorfa} em
$\Omega$ menos um conjunto de polos é
\emph{meromorfa}\index{função meromorfa} em $\Omega$.
\end{definition}

\begin{theorem}[Riemann; Casorati--Weierstrass]
\label{thm:b3:residues:riemanncw}
Seja $f$ \omterm{def:b3:holomorphic:holo}{holomorfa} em $D(a,R)\setminus\{a\}$.
\begin{enumerate}
  \item (Riemann) Se $f$ é \emph{limitada} perto de $a$, a
        singularidade é removível.\index{singularidade removível}
  \item (Casorati--Weierstrass) Se $a$ é essencial, então $f\bigl(D(a,\varepsilon)\setminus\{a\}\bigr)$ é
        denso em $\C$ para todo $\varepsilon$.\index{teorema de
        Casorati--Weierstrass}
\end{enumerate}
\end{theorem}

\begin{proof}
(1) Para $n < 0$ e $\rho \to 0$: $\abs{c_n} \leq
M\rho^{-n-1}\cdot\rho\cdot\rho^{-\,n}\dots$ por ML em $C_\rho$: $\abs{c_n} \leq \frac{1}{2\pi}\,2\pi\rho\cdot
M\rho^{-(n+1)} = M\rho^{-n} \to 0$ (pois
$-n > 0$): todos os coeficientes negativos se anulam.
(2) Se algum valor $b$ não fosse aproximado: $\abs{f - b} \geq
\delta$ perto de
$a$, de modo que $g = 1/(f - b)$ é \omterm{def:b3:holomorphic:holo}{holomorfa} e limitada perto de $a$:
removível (1), e $g$ se estende com valor $c$. Se $c \neq 0$, $f = b + 1/g$ é
limitada perto de $a$: removível --- excluído. Se $c = 0$,
$g$ tem um zero de ordem finita $m$ em $a$
(\cref{thm:b3:holomorphic:identity}; $g \not\equiv 0$), e $f = b + 1/g$ tem um polo de ordem $m$: excluído
de novo.
\end{proof}

\section{O teorema dos resíduos}

\begin{theorem}[Teorema dos resíduos]\label{thm:b3:residues:residue}
Sejam $\Omega$ aberto, $S \subseteq \Omega$ finito, $f \in
\mathcal H(\Omega\setminus S)$ e $\Gamma$ um ciclo em
$\Omega\setminus S$ com $\operatorname{Ind}_\Gamma(w) = 0$ para todo $w \notin \Omega$. Então\index{teorema dos resíduos}
\[
\frac{1}{2\iu\pi}\int_\Gamma f(z)\,\dd z
= \sum_{a\in S}\operatorname{Ind}_\Gamma(a)\,
\operatorname{Res}(f, a) .
\]
\end{theorem}

\begin{proof}
Para cada $a \in S$, seja $P_a(z) = \sum_{n\leq-1}c_n^{(a)}(z -
a)^n$ a parte principal de $f$ em $a$: uma
série que converge em $\C\setminus\{a\}$ (seu raio em $1/(z-a)$ é infinito: a cauda de
Laurent converge para todo $\abs{z-a}$ pequeno, logo, sendo série de potências
em $(z-a)^{-1}$, converge em toda parte) e é ali \omterm{def:b3:holomorphic:holo}{holomorfa}. Então $g = f - \sum_{a\in S}P_a$ tem
singularidades removíveis em cada ponto de $S$ (sua expansão de
Laurent em $a$ não tem parte negativa: as demais $P_{a'}$ são
\omterm{def:b3:holomorphic:holo}{holomorfas} em $a$), de modo que $g$ se estende \omterm{def:b3:holomorphic:holo}{holomorficamente} a
$\Omega$, e o~\cref{thm:b3:residues:globalcauchy} dá $\int_\Gamma g = 0$. Resta integrar cada
$P_a$: termo a termo (convergência normal no \omterm{def:b3:topology:compact}{compacto} $\operatorname{im}\Gamma$, que
evita $a$),
\[
\frac1{2\iu\pi}\int_\Gamma(z - a)^n\,\dd z = 0 \ (n \leq -2:
\text{primitiva } \tfrac{(z-a)^{n+1}}{n+1}),
\qquad
\frac1{2\iu\pi}\int_\Gamma\frac{\dd z}{z - a} =
\operatorname{Ind}_\Gamma(a),
\]
de modo que $\frac1{2\iu\pi}\int_\Gamma P_a =
c_{-1}^{(a)}\operatorname{Ind}_\Gamma(a)$. Some sobre $a$.
\end{proof}

\begin{method}[Calculando resíduos]\label{met:b3:residues:compute}
Polo simples: $\operatorname{Res}(f, a) = \lim_{z\to a}(z -
a)f(z)$; para $f = g/h$ com $g(a) \neq 0$, $h(a) = 0$, $h'(a)
\neq 0$: $\operatorname{Res} = g(a)/h'(a)$. Polo de
ordem $m$: $\operatorname{Res}(f, a) = \frac1{(m-1)!}\lim_{z\to
a}\bigl((z-a)^mf(z)\bigr)^{(m-1)}$. Singularidades essenciais: expanda e leia
$c_{-1}$ (por exemplo, a partir de séries conhecidas). Verifique
sempre quais polos o contorno de fato circunda, e com que índice.
\end{method}

\begin{example}[Os quatro tipos clássicos de integral]
\label{ex:b3:residues:integrals}
(a) \emph{Racional sobre $\R$}: para $\int_\R\frac{\dd x}{1 +
x^4}$, feche com um semicírculo
grande $S_R$ no semiplano superior: o integrando é ali $O(R^{-4})$, de modo
que $\int_{S_R} \to 0$ (ML), e o teorema dos \omterm{def:b3:residues:singularities}{resíduos}, com os polos $\eu^{\iu\pi/4}, \eu^{3\iu\pi/4}$ (simples,
de \omterm{def:b3:residues:singularities}{resíduos} $\frac1{4z^3} = \frac{z}{4z^4} = -\frac z4$ em um polo), dá
\[
\int_\R\frac{\dd x}{1 + x^4}
= 2\iu\pi\Bigl(-\frac{\eu^{\iu\pi/4}}4 -
\frac{\eu^{3\iu\pi/4}}4\Bigr)
= \frac{\pi}{\sqrt2} .
\]
(b) \emph{Tipo Fourier}: para $t \geq 0$, $\int_\R\frac{\eu^{\iu tx}}{1 + x^2}\dd x =
2\iu\pi\operatorname{Res}\bigl(\tfrac{\eu^{\iu tz}}{1+z^2},
\iu\bigr) = 2\iu\pi\frac{\eu^{-t}}{2\iu} = \pi\eu^{-t}$ --- o semicírculo superior
funciona porque $\abs{\eu^{\iu tz}} =
\eu^{-t\operatorname{Im}z} \leq 1$ ali; tomando partes reais: $\int_\R\frac{\cos(tx)}{1+x^2}\dd x = \pi\eu^{-\abs t}$, o que liquida a
fórmula admitida no~\cref{exo:b3:lebesgue:10}.
(c) \emph{Trigonométrica sobre um período}: substitua $z =
\eu^{\iu t}$, $\cos t = \frac{z + z^{-1}}2$,
$\dd t =
\frac{\dd z}{\iu z}$: $\int_0^{2\pi}\frac{\dd t}{a + \cos t}$ ($a > 1$) torna-se uma contagem de \omterm{def:b3:residues:singularities}{resíduos} dentro do
círculo unitário (\cref{exo:b3:residues:2}).
(d) \emph{Séries}: emparelhe $f$ com $\pi\cot(\pi z)$, cujos polos são os
inteiros com \omterm{def:b3:residues:singularities}{resíduo} $1$: o problema de fim de semana soma $\sum n^{-2}$ e
$\sum n^{-4}$ por esse \omterm{def:b3:topology:pathconnected}{caminho}.
\end{example}

\begin{omfigure}
\begin{tikzpicture}[scale=1.15]
  \draw[->] (-2.5,0) -- (2.6,0) node[right, font=\small]
    {$\operatorname{Re}$};
  \draw[->] (0,-0.5) -- (0,2.5) node[above, font=\small]
    {$\operatorname{Im}$};
  \draw[omThm, very thick, ->] (-2,0) -- (0.3,0);
  \draw[omThm, very thick] (0.3,0) -- (2,0);
  \draw[omThm, very thick] (2,0) arc (0:90:2);
  \draw[omThm, very thick, ->] (2,0) arc (0:50:2);
  \draw[omThm, very thick] (0,2) arc (90:180:2);
  \node[omThm, font=\small] at (1.75,1.55) {$S_R$};
  \node[font=\small] at (2.2,-0.25) {$R$};
  \node[font=\small] at (-2.15,-0.25) {$-R$};
  \fill[omDef] (45:1.19) circle (1.6pt)
    node[above right, font=\small] {$\eu^{\iu\pi/4}$};
  \fill[omDef] (135:1.19) circle (1.6pt)
    node[above left, font=\small] {$\eu^{3\iu\pi/4}$};
  \fill[black!40] (-45:1.19) circle (1.6pt);
  \fill[black!40] (-135:1.19) circle (1.6pt);
\end{tikzpicture}

{\small O contorno semicircular para $\int_\R\frac{\dd x}{1 +
x^4}$: quando $R \to \infty$, o arco
contribui com $O(R^{-3})$, e o teorema dos \omterm{def:b3:residues:singularities}{resíduos} conta os dois polos
englobados (em azul). Os dois polos inferiores (em cinza) ficam de
fora: índice $0$.}
\end{omfigure}

\section{O princípio do argumento e o teorema de Rouché}

\begin{theorem}[Princípio do argumento]
\label{thm:b3:residues:argument}
Sejam $f$ \omterm{def:b3:residues:singularities}{meromorfa} em $\Omega$, com zeros $z_j$ (de ordens
$m_j$) e polos $p_k$ (de ordens $\mu_k$), e $\gamma$ um
\omterm{def:b3:topology:pathconnected}{caminho} fechado em $\Omega$ que os evita todos, com $\operatorname{Ind}_\gamma = 0$ fora de
$\Omega$. Então\index{princípio do argumento}
\[
\frac1{2\iu\pi}\int_\gamma\frac{f'(z)}{f(z)}\,\dd z
= \sum_j m_j\operatorname{Ind}_\gamma(z_j)
- \sum_k \mu_k\operatorname{Ind}_\gamma(p_k)
\]
(apenas um número finito de termos é não nulo). Para um contorno simples
percorrido no sentido anti-horário, a integral conta os zeros menos os
polos \omterm{def:b3:topology:interior}{interiores}, com multiplicidade --- e é igual ao índice de rotação
do \omterm{def:b3:topology:pathconnected}{caminho} imagem $f\circ\gamma$ em torno de $0$.
\end{theorem}

\begin{proof}
Perto de um zero de ordem $m$: $f = (z-a)^mg$, $g(a) \neq 0$, de modo que $\frac{f'}f = \frac m{z - a} + \frac{g'}g$ com o
segundo termo \omterm{def:b3:holomorphic:holo}{holomorfo} perto de $a$: um polo simples de \omterm{def:b3:residues:singularities}{resíduo}
$m$. Perto de um polo de ordem $\mu$: $f = (z-a)^{-\mu}g$ dá \omterm{def:b3:residues:singularities}{resíduo}
$-\mu$. Nos demais pontos, $\frac{f'}f$ é \omterm{def:b3:holomorphic:holo}{holomorfa}. (Os zeros e polos de
índice não nulo estão em uma região compacta circundada por
$\gamma$; pelo teorema da identidade, ali são em número finito,
$f \not\equiv 0$.) Aplique o~\cref{thm:b3:residues:residue}. A última observação: $\frac1{2\iu\pi}\int_\gamma\frac{f'}f =
\frac1{2\iu\pi}\int_{f\circ\gamma}\frac{\dd w}w =
\operatorname{Ind}_{f\circ\gamma}(0)$ (substitua
$w =
f(\gamma(t))$).
\end{proof}

\begin{theorem}[Rouché]\label{thm:b3:residues:rouche}
Sejam $f, g$ \omterm{def:b3:holomorphic:holo}{holomorfas} em $\Omega$ e $\gamma$ um \omterm{def:b3:topology:pathconnected}{caminho}
fechado com $\operatorname{Ind}_\gamma \in \{0,1\}$, nulo fora de $\Omega$ (um contorno simples).
Se\index{teorema de Rouché}
\[
\abs{g(z)} < \abs{f(z)} \qquad \text{em }
\operatorname{im}\gamma,
\]
então $f$ e $f + g$ têm o mesmo número de zeros (com
multiplicidade) na região $\{\operatorname{Ind}_\gamma =
1\}$.
\end{theorem}

\begin{proof}
Para $t \in \intcc01$, $f_t = f + tg$ não tem zero em $\operatorname{im}\gamma$ ($\abs{f_t} \geq \abs f - \abs g >
0$), de modo que
\[
N(t) = \frac1{2\iu\pi}\int_\gamma
\frac{f_t'(z)}{f_t(z)}\,\dd z
\]
está bem definido; ele conta os zeros na região englobada
(\cref{thm:b3:residues:argument}; não há polos). $N$ é \omterm{def:b3:topology:continuity}{contínuo} em $t$ (o
integrando é conjuntamente \omterm{def:b3:topology:continuity}{contínuo} e os denominadores são
uniformemente limitados por baixo --- convergência dominada) e assume
valores inteiros: é constante. $N(0) = N(1)$.
\end{proof}

\begin{corollary}[Teorema da aplicação aberta]
\label{cor:b3:residues:openmapping}
Uma \omterm{def:b3:holomorphic:holo}{função holomorfa} não constante em um aberto \omterm{def:b3:topology:connected}{conexo} é uma aplicação
aberta. Em particular (de novo), vale o princípio do máximo, e uma
bijeção \omterm{def:b3:holomorphic:holo}{holomorfa} tem inversa \omterm{def:b3:holomorphic:holo}{holomorfa}.\index{teorema da aplicação
aberta (holomorfa)}
\end{corollary}

\begin{proof}
Seja $f(a) = b$; $f - b$ tem um zero de ordem finita $m
\geq 1$ em $a$
(teorema da identidade: $f \not\equiv b$). Escolha $r$ com $f - b$ sem zeros
em $\bar D(a, r)\setminus\{a\}$ (\omterm{thm:b3:holomorphic:identity}{zeros isolados}) e ponha $\delta = \min_{\abs{z - a} =
r}\abs{f(z) - b} > 0$. Para $\abs{w - b} < \delta$: no círculo, $\abs{(b - w)} < \delta \leq \abs{f - b}$,
de modo que Rouché ($f - b$ contra a constante $b - w$) diz que
$f - w$ tem exatamente $m$ zeros em $D(a, r)$: todo tal $w$ é
atingido --- $f(D(a,r)) \supseteq D(b, \delta)$: é aberta. Princípio do máximo: um máximo \omterm{def:b3:topology:interior}{interior} de
$\abs f$ é impossível para $f$ não constante, pois sua imagem em
torno de $f(a)$ contém pontos de módulo maior. Inversa: uma bijeção
\omterm{def:b3:holomorphic:holo}{holomorfa} $f$ é aberta, de modo que $f^{-1}$ é \omterm{def:b3:topology:continuity}{contínua}; o zero
de $f - f(a)$ em $a$ é simples ($m \geq 2$ daria $m$ pré-imagens de valores
próximos --- distintas, pois $f'$ só se anula em pontos isolados, de
modo que perto de $a$ os $m$ zeros de $f - w$ são simples e
distintos para $w$ pequeno genérico: contradição com a
injetividade); então $f'(a) \neq 0$ e o quociente de diferenças de $f^{-1}$
converge: $\bigl(f^{-1}\bigr)'(b) =
1/f'(a)$.
\end{proof}

\section{Exercícios}

\begin{exercise}[$\star$]\label{exo:b3:residues:1}
Classifique a singularidade em $0$ e calcule o \omterm{def:b3:residues:singularities}{resíduo}:
\[
\frac{\sin z}{z},\qquad
\frac{\eu^z - 1}{z^2},\qquad
\frac{1}{z(z-1)^2},\qquad
\frac{\cos z}{z^3},\qquad
\eu^{1/z},\qquad
\frac1{\sin z} .
\]
Dê também o \omterm{def:b3:residues:singularities}{resíduo} da terceira em $z = 1$ e o da última em $z = \pi$.
\end{exercise}

\begin{exercise}[$\star$]\label{exo:b3:residues:2}
Para $a > 1$, calcule, via $z = \eu^{\iu t}$:
\[
\int_0^{2\pi}\frac{\dd t}{a + \cos t}
= \frac{2\pi}{\sqrt{a^2 - 1}} .
\]
Confira os comportamentos-limite $a \to 1^+$ e $a \to \infty$.
\end{exercise}

\begin{exercise}[$\star\star$]\label{exo:b3:residues:3}
Calcule com contornos semicirculares, justificando as estimativas nos
arcos:
\[
\int_\R\frac{x^2}{1 + x^6}\,\dd x = \frac\pi3,
\qquad
\int_\R\frac{\dd x}{(1 + x^2)^{2}} = \frac\pi2
\quad\text{(reobtendo \cref{exo:b3:fouriertransform:3})}.
\]
\end{exercise}

\begin{exercise}[$\star\star$]\label{exo:b3:residues:4}
Demonstre, para $t \geq 0$ e $a > 0$:
\[
\int_\R\frac{\cos(tx)}{x^2 + a^2}\,\dd x =
\frac{\pi}{a}\,\eu^{-at},
\]
e deduza a transformada de Fourier de $x \mapsto \eu^{-a\abs
x}$ por inversão --- comparando
com o~\cref{exo:b3:fouriertransform:1}.
\end{exercise}

\begin{exercise}[$\star\star$]\label{exo:b3:residues:5}
Expanda $f(z) = \dfrac1{(z-1)(z-2)}$ em \omterm{thm:b3:residues:laurent}{série de Laurent} em cada uma das três regiões $\abs z < 1$,
$1 < \abs z < 2$, $\abs z >
2$. Por que as três expansões diferem? Explique por que o
coeficiente de $z^{-1}$ na segunda e na terceira expansões
\emph{não} é um \omterm{def:b3:residues:singularities}{resíduo} de $f$ em $0$ ($f$ não tem
singularidade ali) e calcule os \omterm{def:b3:residues:singularities}{resíduos} de fato de $f$, em $1$
e em $2$.
\end{exercise}

\begin{exercise}[$\star\star$]\label{exo:b3:residues:6}
(a) Mostre que $\eu^{1/z}$ tem singularidade essencial em $0$ e verifique
Casorati--Weierstrass à mão: resolva $\eu^{1/z} =
w$ explicitamente para qualquer
$w \neq 0$, exibindo soluções arbitrariamente próximas de $0$.
(b) Mostre que $\abs{\eu^{1/z}}$ é ilimitada em toda \omterm{def:b3:topology:topology}{vizinhança} perfurada de
$0$ e, ainda assim, $\eu^{1/z}$ não tem polo: que limite falha?
\end{exercise}

\begin{exercise}[$\star\star$]\label{exo:b3:residues:7}
Conte com Rouché:
(a) os zeros de $z^7 - 4z^3 + z - 1$ em $\abs z < 1$;
(b) os zeros de $z^4 + 5z + 1$ em $\abs z < 1$ e em $1 <
\abs z < 2$;
(c) redemonstre d'Alembert--Gauss: um polinômio mônico de grau $n$
tem $n$ zeros em algum disco grande \emph{(compare com $z^n$)}.
\end{exercise}

\begin{exercise}[$\star\star\star$]\label{exo:b3:residues:8}
(Hurwitz) Sejam $f_n \to f$ uniformemente nos \omterm{def:b3:topology:compact}{compactos}, $f_n \in
\mathcal H(\Omega)$, $\Omega$
\omterm{def:b3:topology:connected}{conexo}, $f \not\equiv 0$.
(a) Mostre que, se todas as $f_n$ são sem zeros, $f$ também é.
\emph{(Se $f(a) = 0$: princípio do argumento em um pequeno círculo em torno
de $a$, e o~\cref{thm:b3:holomorphic:weierstrassconv} para passar ao limite em $\int f_n'/f_n$.)}
(b) Mostre que, se todas as $f_n$ são injetoras, então $f$ é
injetora ou constante. \emph{(Aplique (a) a $z \mapsto f_n(z) - f_n(w)$ em $\Omega\setminus\{w\}$.)}
\end{exercise}

\begin{exercise}[$\star\star\star$]\label{exo:b3:residues:9}
Para $n \geq 2$, integre $\frac1{1 + z^n}$ na fronteira do setor $\{0 \leq \arg z \leq \frac{2\pi}n,\ \abs z
\leq R\}$ e deduza
\[
\int_0^{+\infty}\frac{\dd x}{1 + x^n} =
\frac{\pi}{n\,\sin(\pi/n)} .
\]
Verifique $n = 2$ contra $\arctan$, e o limite $n \to
\infty$.
\end{exercise}

\begin{exercise}[$\star\star$]\label{exo:b3:residues:10}
Seja $f$ uma função racional com $\deg(\text{denominador})
\geq \deg(\text{numerador}) + 2$. Mostre que a soma de
\emph{todos} os \omterm{def:b3:residues:singularities}{resíduos} de $f$ é zero \emph{(integre sobre círculos
cada vez maiores)}. Use isso para recalcular a decomposição em frações
parciais de $\frac1{z(z-1)(z-2)}$ sem nenhuma álgebra linear.
\end{exercise}

\begin{exercise}[$\star\star\star$]\label{exo:b3:residues:11}
(O buraco de fechadura: a integral de reflexão de Euler) Para $0 < a < 1$,
calcule
\[
I(a) = \int_0^{\infty}\frac{x^{a-1}}{1 + x}\,\dd x =
\frac{\pi}{\sin(\pi a)}
\]
integrando $f(z) = \frac{z^{a-1}}{1+z} =
\frac{\eu^{(a-1)\log z}}{1 + z}$ (logaritmo cortado ao longo de $\R_+$, $\arg z \in \intoo0{2\pi}$) sobre
o contorno em buraco de fechadura: para fora ao longo da margem
superior do corte, de $\varepsilon$ a $R$, em torno de $C_R$, de volta
sob o corte, em torno de $C_\varepsilon$. Justifique: os dois trechos retos
diferem pelo fator $\eu^{2\iu\pi(a-1)}$, as contribuições dos círculos se anulam
($R^{a-1}\cdot R \to 0$ e $\varepsilon^{a-1}\cdot\varepsilon \to 0$), e o único polo $z = -1$ tem \omterm{def:b3:residues:singularities}{resíduo} $\eu^{\iu\pi(a - 1)}$. Deduza
também $\Gamma(a)\Gamma(1 - a) = \frac\pi{\sin\pi a}$ \emph{(escreva $\Gamma(a)\Gamma(1-a) = B(a, 1-a)$ pelo~\cref{pb:b3:lebesgue:1} e substitua $t =
\frac{x}{1+x}$)}.
\end{exercise}

\begin{exercise}[$\star\star$]\label{exo:b3:residues:12}
(Contando zeros com o princípio do argumento, numericamente) Seja
$P(z) = z^4 + 8z + 1$.
(a) Quantos zeros no disco unitário? \emph{(Rouché contra
$8z + 1$.)}
(b) Quantos no anel $1 < \abs z < 3$? \emph{(Rouché contra $z^4$ em $\abs z = 3$.)}
Refine: mostre que todo zero tem módulo $< 2.1$.
(c) Quantos no semiplano direito? \emph{(Conte primeiro em $\abs z = 2$; depois
acompanhe a imagem do eixo imaginário: $P(\iu t) = t^4 + 1 + 8\iu t$ tem parte real positiva em
todo ele, de modo que não há zeros no eixo, e a variação do argumento
ao longo dele é calculável --- conclua com um semidisco grande.)}
\end{exercise}

\section{Problema: \texorpdfstring{$\zeta(2k)$}{zeta(2k)} pela cotangente}

\begin{problem}[{Problema de fim de semana --- somando $\sum n^{-2k}$ com
\omterm{def:b3:residues:singularities}{resíduos}}]\label{pb:b3:residues:1}
O teorema dos \omterm{def:b3:residues:singularities}{resíduos} soma séries: emparelhar uma função racional com
$\pi\cot(\pi z)$, cujos polos ficam nos inteiros, transforma $\sum_{n}f(n)$ em uma contagem
de \omterm{def:b3:residues:singularities}{resíduos}. Demonstramos o método e calculamos $\zeta(2) = \frac{\pi^2}{6}$ e $\zeta(4) =
\frac{\pi^4}{90}$ --- os
valores encontrados por séries de Fourier no segundo ano e por traços de
operadores no~\cref{ch:b3:spectral}, agora por integração de contorno.

\medskip
\noindent\textbf{Parte I --- O núcleo cotangente.}
\begin{enumerate}
  \item Mostre que $\pi\cot(\pi z)$ é \omterm{def:b3:residues:singularities}{meromorfa} em $\C$, com polos simples
        exatamente em $z = n \in \Z$, cada um de \omterm{def:b3:residues:singularities}{resíduo} $1$ \emph{(calcule
        $\lim_{z\to n}(z-n)\pi\cot\pi z$)}.
  \item Calcule o início da expansão de Laurent em $0$:
        \[
        \pi\cot(\pi z) = \frac1z - \frac{\pi^2}{3}\,z -
        \frac{\pi^4}{45}\,z^3 + O(z^5) ,
        \]
        dividindo a série de potências de $\cos$ pela de $\sin$
        (justifique a divisão: $\frac{\sin\pi z}{\pi z}$ é \omterm{def:b3:holomorphic:holo}{holomorfa} e não nula perto de
        $0$, de modo que seu inverso é \omterm{def:b3:holomorphic:holo}{holomorfo}; identifique os
        coeficientes até a ordem $3$).
  \item Seja $C_N$ a fronteira do quadrado de vértices $(\pm1\pm\iu)(N + \frac12)$. Mostre
        que $\abs{\cot(\pi z)} \leq 2$ em $C_N$ para todo $N \geq
        1$. \emph{(Nos lados
        verticais, $\cot(\pi(\pm(N +
        \frac12) + \iu y)) = \mp\tan(\iu\pi y)$, de módulo $\abs{\tanh(\pi y)} \leq 1$; nos lados horizontais $\abs y = N + \frac12$,
        limite $\abs{\cot(\pi(x\pm\iu y))} \leq \coth(\pi y) \leq
        \coth(\pi/2) < 1.1$.)}
\end{enumerate}

\medskip
\noindent\textbf{Parte II --- O teorema de somação.}
\begin{enumerate}[resume]
  \item Seja $f$ racional, \omterm{def:b3:holomorphic:holo}{holomorfa} nos inteiros, com $\deg(\text{denom}) \geq \deg(\text{numer}) + 2$. Usando
        o teorema dos \omterm{def:b3:residues:singularities}{resíduos} em $C_N$ e a cota da questão 3,
        demonstre:
        \[
        \lim_{N\to\infty}\ \sum_{n = -N}^{N} f(n)
        = -\sum_{p\ \text{polo de}\ f}
        \operatorname{Res}\bigl(\pi\cot(\pi z)f(z),\,p\bigr).
        \]
  \item Onde o argumento precisa da condição sobre o grau? Mostre com
        um exemplo (tome $f(z) = 1/(z + \frac12)$) que, para decaimento mais lento, o
        limite simétrico ainda pode existir enquanto a série de duas
        pontas diverge --- e que a fórmula calcula então o \emph{valor
        principal}.
\end{enumerate}

\medskip
\noindent\textbf{Parte III --- Os valores.}
\begin{enumerate}[resume]
  \item Aplique o método a $f(z) = 1/z^2$: aqui $f$ tem seu polo \emph{em} um
        inteiro, de modo que se deve rodar o argumento diretamente ---
        integre $g(z) = \frac{\pi\cot(\pi
        z)}{z^2}$ sobre $C_N$, mostre que a integral $\to 0$
        e calcule $\operatorname{Res}(g, 0)$ a partir da questão 2. Conclua:
        \[
        2\sum_{n\geq1}\frac1{n^2} = \frac{\pi^2}{3},
        \qquad \zeta(2) = \frac{\pi^2}6 .
        \]
  \item O mesmo com $g(z) = \frac{\pi\cot(\pi z)}{z^4}$: calcule $\operatorname{Res}(g, 0)$ e deduza $\zeta(4) = \frac{\pi^4}{90}$.
  \item Explique o padrão geral: para todo $k \geq 1$, $\zeta(2k)$ é $-\frac12$ vezes o
        coeficiente de $z^{2k-1}$ na expansão de Laurent de $\pi\cot(\pi
        z)$ em $0$
        --- um múltiplo racional de $\pi^{2k}$. Calcule $\zeta(6)$ levando a
        divisão da questão 2 um passo adiante. O que o método diz sobre
        $\zeta(3)$ --- e por que ele não diz nada?
\end{enumerate}

\medskip
\noindent\textbf{Parte IV --- A expansão em frações parciais da
cotangente.}
\begin{enumerate}[resume]
  \item Fixe $w \in \C\setminus\Z$ e aplique o método da Parte II a $f(z) = \dfrac{1}{(z - w)(z + w)}$ --- notando
        que $\pi\cot(\pi z)f(z)$ tem agora polos simples adicionais em $\pm w$,
        cujos \omterm{def:b3:residues:singularities}{resíduos} devem entrar na conta. Deduza a expansão em
        frações parciais
        \[
        \pi\cot(\pi w) = \frac1w +
        \sum_{n\geq1}\frac{2w}{w^2 - n^2},
        \]
        convergindo a série normalmente nos \omterm{def:b3:topology:compact}{compactos} de $\C\setminus\Z$.
  \item Recupere dessa expansão, expandindo cada termo em potências de
        $w$ (justifique a troca), os \emph{mesmos} coeficientes de
        Laurent da questão 2 --- o círculo se fecha: a fórmula de Euler
        $\sum\frac1{n^2} = \frac{\pi^2}6$ é o coeficiente de $w$ nas duas faces da cotangente.
        Compare com a demonstração por séries de Fourier (segundo ano)
        e com a demonstração por traço (\cref{pb:b3:spectral:1}): três teorias, um
        número.
\end{enumerate}

\medskip
\noindent\textbf{Parte V --- O produto de Euler para o seno.} A
expansão da questão 9 é a derivada logarítmica de um produto infinito;
demonstramos agora honestamente a fatoração de Euler de 1734.
\begin{enumerate}[resume]
  \item Para $N \geq 1$, ponha $P_N(z) =
        z\prod_{n=1}^{N}\bigl(1 - \frac{z^2}{n^2}\bigr)$. Mostre que $P_N$ converge,
        uniformemente em todo disco $\bar D(0, R)$, para uma função inteira
        $P$ cujos zeros são exatamente os inteiros, todos simples.
        \emph{(Para $n
        \geq 2R$, escreva o fator como $\exp\log(1 -
        z^2/n^2)$ com o logaritmo
        principal do~\cref{exo:b3:holomorphic:3}, limite $\abs{\log(1+u)}
        \leq 2\abs u$ para $\abs u \leq \frac12$ pela série
        e exponencie a soma normalmente convergente de logaritmos; os
        finitos fatores restantes formam um polinômio. Conclua com
        o~\cref{thm:b3:holomorphic:weierstrassconv}.)}
  \item Mostre que, em $\C\setminus\Z$,
        \[
        \frac{P'(z)}{P(z)} = \frac1z +
        \sum_{n\geq1}\frac{2z}{z^2 - n^2} = \pi\cot(\pi z)
        \]
        \emph{(derive os produtos finitos, passe ao limite usando
        o~\cref{thm:b3:holomorphic:weierstrassconv} e a ausência de zeros de $P$ fora de
        $\Z$, e cite a questão 9)}.
  \item Mostre que $Q = \sin(\pi z)/P(z)$ se estende a uma função inteira sem zeros com
        $Q' = 0$, e conclua o \emph{produto de Euler}:
        \[
        \sin(\pi z) = \pi z\prod_{n\geq1}
        \Bigl(1 - \frac{z^2}{n^2}\Bigr)
        \qquad (z \in \C) .
        \]
  \item (Wallis, 1655) Avalie em $z = \frac12$:
        \[
        \frac\pi2 = \prod_{n\geq1}\frac{4n^2}{4n^2 - 1}
        = \lim_{N\to\infty}
        \frac{2\cdot2\cdot4\cdot4\cdots(2N)(2N)}
        {1\cdot3\cdot3\cdot5\cdots(2N-1)(2N+1)} .
        \]
  \item Para $\abs z < 1$, expanda o logaritmo do produto como série dupla
        (justifique o rearranjo) e recupere $\zeta(2) =
        \frac{\pi^2}6$, comparando o
        coeficiente de $z^3$ em $\sin(\pi z) = \pi z - \frac{\pi^3}6z^3 + \cdots$ --- a face do produto do número
        de Euler.
\end{enumerate}

\medskip
\noindent\textbf{Parte VI --- Núcleos irmãos.} A cotangente tem irmãos;
cada um precifica sua própria família de séries.
\begin{enumerate}[resume]
  \item Derive termo a termo a expansão da questão 9 (justificado
        pelo~\cref{thm:b3:holomorphic:weierstrassconv}) para obter, normalmente nos \omterm{def:b3:topology:compact}{compactos} de
        $\C\setminus\Z$,
        \[
        \frac{\pi^2}{\sin^2(\pi z)} =
        \sum_{n\in\Z}\frac1{(z - n)^2} .
        \]
  \item Avalie em $z = \frac12$: $\sum_{m\geq0}\frac1{(2m+1)^2} = \frac{\pi^2}8$; recupere $\zeta(2)$ mais uma vez, separando os
        inteiros por paridade.
  \item Verifique a identidade de duplicação $\tan\theta =
        \cot\theta - 2\cot(2\theta)$ e deduza
        \[
        \pi\tan(\pi z) =
        \sum_{m\geq0}\frac{8z}{(2m+1)^2 - 4z^2} ,
        \]
        normalmente nos \omterm{def:b3:topology:compact}{compactos} que evitam $\frac12 + \Z$.
  \item Expanda em torno de $0$ ($\abs z < \frac12$; Fubini de novo): com $\lambda(s) =
        \sum_{m\geq0}(2m+1)^{-s}$,
        \[
        \pi\tan(\pi z) =
        \sum_{k\geq0}8\cdot4^k\,\lambda(2k+2)\,z^{2k+1} ;
        \]
        compare com $\tan u = u + \frac{u^3}3 + O(u^5)$ para recuperar $\lambda(2) = \frac{\pi^2}8$ e obter $\lambda(4) = \frac{\pi^4}{96}$, e confira
        então $\zeta(4) = \frac{\pi^4}{90}$ via $\lambda(4) = (1 -
        2^{-4})\,\zeta(4)$.
  \item Verifique $\frac1{\sin\theta} = \cot\frac\theta2 -
        \cot\theta$ e deduza
        \[
        \frac{\pi}{\sin(\pi z)} = \frac1z +
        \sum_{n\geq1}(-1)^n\,\frac{2z}{z^2 - n^2} .
        \]
        Confira os sinais contra os \omterm{def:b3:residues:singularities}{resíduos} de $\pi/\sin(\pi z)$ nos inteiros.
  \item Leia o coeficiente de $z$: $\eta(2) =
        \sum_{n\geq1}\frac{(-1)^{n-1}}{n^2} =
        \frac{\pi^2}{12}$, e confirme a coerência
        $\eta(2) = (1 - 2^{1-2})\,\zeta(2)$.
  \item (Final) Avalie a expansão da questão 20 em $z =
        \frac12$ e deduza a
        \emph{fórmula de Leibniz}
        \[
        \frac\pi4 = 1 - \frac13 + \frac15 - \frac17 +
        \cdots
        \]
        Encerre com um parágrafo curto: um núcleo por aritmética ---
        que núcleo precifica que família de séries, e por que todos eles
        são estruturalmente cegos a $\zeta(3)$.

\end{enumerate}
\medskip
\noindent\textbf{Parte VII --- A tabela de preços \omterm{def:b3:complete:complete}{completa}: os números
de Bernoulli.}
\begin{enumerate}[resume]
  \item Combine a expansão em frações parciais de $\pi
        z\cot(\pi z)$ com a função
        geradora dos números de Bernoulli ($\frac{w}{\eu^w - 1} =
        \sum_n\frac{B_n}{n!}w^n$, \cref{pb:b3:holomorphic:1}, Parte
        VI): a partir de
        \[
        \pi z\cot(\pi z) = \iu\pi z +
        \frac{2\iu\pi z}{\eu^{2\iu\pi z} - 1}
        \]
        (demonstre primeiro essa identidade), deduza a forma fechada
        \[
        \zeta(2k) = (-1)^{k+1}\,
        \frac{(2\pi)^{2k}\,B_{2k}}{2\,(2k)!}
        \qquad (k \geq 1).
        \]
  \item Confira a fórmula em $B_2 = \frac16$, $B_4 =
        -\frac1{30}$, $B_6 = \frac1{42}$: recupere $\zeta(2) =
        \frac{\pi^2}6$, $\zeta(4) = \frac{\pi^4}{90}$ e
        calcule $\zeta(6) = \frac{\pi^6}{945}$.
  \item (Recursão de Euler) Expanda os dois lados de $\bigl(z\cot z\bigr)' = \cot z - z(1 + \cot^2z)$ --- ou eleve
        ao quadrado diretamente a série da cotangente --- para
        demonstrar
        \[
        \Bigl(k + \frac12\Bigr)\zeta(2k) =
        \sum_{j=1}^{k-1}\zeta(2j)\,\zeta(2k - 2j)
        \qquad (k \geq 2),
        \]
        e verifique que ela calcula $\zeta(4)$ a partir de $\zeta(2)$ e $\zeta(6)$ a
        partir de $\zeta(2), \zeta(4)$ --- todos os valores zeta pares a partir da
        única semente $\frac{\pi^2}6$, sem nenhuma integração nova.
\end{enumerate}
\end{problem}
